\documentclass[12pt]{article}
\usepackage[utf8]{inputenc}
\usepackage[T1]{fontenc}
\usepackage{hyperref}
\title{Tulil Language Resources for K-12 Curriculum Development}
\author{Compiled by Claude Sonnet 5 (claude-sonnet-5) for Andr\'as Kornai}
\date{}
\begin{document}
\maketitle
\section{Introduction}
Tulil (Taulil, ISO 639-3 tuh) is a Papuan language of East New Britain Province, Papua New Guinea. It belongs to the small Taulil--Butam family; Butam is extinct. Sources in my search give about 2,000 speakers and list the language as vulnerable. Very little educational material was found.
\section{Primary Educational Resources}
No primers, textbooks or school materials for Tulil turned up in the search results.
\section{Linguistic and Reference Materials}
\begin{itemize}
\item Wikipedia, Taulil language (overview): \url{https://en.wikipedia.org/wiki/Taulil_language}
\item Wikipedia, Taulil--Butam languages: \url{https://en.wikipedia.org/wiki/Taulil%E2%80%93Butam_languages}
\item Endangered Languages Project entry: \url{https://endangeredlanguages.com/node/87565}
\end{itemize}
The search results also mention a PhD thesis, Meng (2018), \emph{A grammar of Tulil} (La Trobe University). I did not see a direct URL for it in the results, so it is not listed as a resource here. No dictionary or wordlist was found.
\section{Religious and Community Resources}
Joshua Project lists Scripture translation as needed and mentions community interest in a planned SIL New Britain Initiative: \url{https://joshuaproject.net/languages/tuh}
\section{Digital and Audio Resources}
Global Recordings Network has a page for the language but reports no recordings available: \url{https://globalrecordings.net/en/language/tuh}
\section{Curriculum Development Context}
The language is not known to be taught in schools. A curriculum would have to start from the Meng grammar and new community-led documentation. Any native-language curriculum would be built largely from scratch.
\end{document}